\documentclass["../Cours.tex"]{subfiles}
\usepackage{etoolbox}

\begin{document}
\chapitre{Fractions et proportionnalité}

\partie{Calculs fractionnaires}

\definition{Soient deux nombres $a$ et $b$, avec $b \neq 0$. Le quotient de $a$ par $b$ est le nombre qui, lorsqu'il est multiplié par $b$, donne $a$. $$\dfrac{a}{b} \times b = a$$}

\notation{On le note $\dfrac{a}{b}$. (écriture fractionnaire)}

\begin{listedexemples}
    \item $\dfrac{3}{4} \times 4 = \dfrac{3}{\cancel{4}} \times \cancel{4} = 3$
    \item $\dfrac{7,2}{5,8} \times 5,8 = \dfrac{7,2}{\cancel{5,8}} \times \cancel{5,8} = 7,2$
\end{listedexemples}


\propriete{Soit deux nombres $a$ et $b$ avec $b \neq 0$. Soit $k$ avec $k \neq 0$. 
$$\dfrac{a}{b} = \dfrac{ka}{kb}$$
}

\begin{listedexemples}
    \item $\dfrac{16}{20} = \dfrac{16 \times 4}{20 \times 4} = \dfrac{64}{80}$
    \item $\dfrac{35}{75} = \dfrac{35 \div 5}{75 \div 5} = \dfrac{7}{15}$
\end{listedexemples}

\definition{Simplifier une fraction, c'est trouver une autre fraction qui lui est égale avec un numérateur et un dénominateur plus petits.}

\propriete{Pour additionner deux fractions, il faut qu'elles aient le même dénominateur.}

\propriete{Soient deux nombres $a$ et $b$, $b \neq 0$ et deux nombres $c$ et $d$, $d \neq 0$. $$ \dfrac{a}{b} + \dfrac{c}{d} = \dfrac{ad+cb}{bd}$$}

\begin{listedexemples}
    \item $\dfrac{3}{7} + \dfrac{12}{8} = \dfrac{3 \times 8}{7 \times 8} + \dfrac{12 \times 7}{8 \times 7}$ 
\end{listedexemples}

\propriete{Soient $a$ et $b$ deux nombres, $b \neq 0$.\\ Soient $c$ et $d$ deux nombres, $d \neq 0$. $$\dfrac{a}{b} \times \dfrac{c}{d} = \dfrac{a \times c}{b \times d}$$}

\exemple{$$\dfrac{7}{5} \times \dfrac{3}{4} = \dfrac{7 \times 3}{5 \times 4} = \dfrac{21}{20}$$}

\propriete{Diviser par un nombre revient à multiplier par son inverse.}

\exemple{
\begin{center}
\begin{tikzpicture}
    \node at (0,0) {$\dfrac{2}{3} \div \dfrac{4}{5} = \dfrac{2}{3} \times \dfrac{5}{4} = \dfrac{10}{12}$};
    \draw[rouge] (0.1,-0.6) rectangle (1.1,0.6);
    \draw[rouge] (-1.85,-0.6) rectangle (-0.85,0.6);
    \draw[rouge,thick,-latex] (-1.35,-0.65) arc (-180:0:1 and 0.5);
    \node at (0,-1.5) {la fraction est inversée};
    \node at (0,-2) {et le $\div$ devient $\times$};
\end{tikzpicture}
\end{center}
}

\clearpage
\partie{Pourcentages}

\souspartie{Le pourcentage d'un nombre}

\definition{Calculer $p~\%$ d'un nombre $n$, c'est multiplier $\dfrac{p}{100}$ par $n$.}

\renewcommand\CancelColor{\color{rouge}}

\begin{listedexemples}
    \item $25~\%$ de $12 = \dfrac{25}{100} \times 12 = \dfrac{\cancel{25} \times 1}{\cancel{25} \times \cancel{4}} \times 3 \times \cancel{4} = 3$
    \item $25~\%$ de $100 = \dfrac{25}{\cancel{100}} \times \cancel{100} = 25$
    \item $75~\%$ de $32 = \dfrac{75}{100} \times 32 = \dfrac{3 \times \cancel{25}}{\cancel{4} \times \cancel{25}} \times 8 \times \cancel{4} = 3 \times 8 = 24$
\end{listedexemples}

\regle{Pour calculer une proportion sous forme de pourcentage, il faut que le dénominateur soit égal à 100.}

\exemple{%
    \begin{center}
    \begin{tikzpicture}
        \definecolor{bleuRF}{RGB}{05, 20, 64}
        \definecolor{blancRF}{RGB}{255, 255, 255}
        \definecolor{rougeRF}{RGB}{236, 25, 32}
        \draw[white, fill=bleuRF] (0,0) rectangle (1,2);
        \draw[white, fill=blancRF] (1,0) rectangle (2,2);
        \draw[white, fill=rougeRF] (2,0) rectangle (3,2);
    \end{tikzpicture}
    \end{center}
Dans le drapeau français, il y a 1 morceau sur 3 qui est bleu. Donc, la proportion de bleu dans le drapeau est de $\dfrac{1}{3}$. Pour la mettre sous forme de pourcentage, je peux multiplier au numérateur et au dénominateur par 100 :
$$\dfrac{1}{3} = \dfrac{1 \times 100}{3 \times 100} = \dfrac{1 \times 100}{3}~\% \approx 33,3 ~\% = \dfrac{33{,}3}{100}$$
}

\clearpage
\souspartie{Augmentation et diminution par un pourcentage}

\regle{\begin{itemize}\vspace{-1em}
    \item Augmenter une quantité de $x~\%$, c'est multiplier cette quantité par $\dfrac{100+x}{100}$.
    \item Diminuer une quantité de $x~\%$, c'est multiplier cette quantité par $\dfrac{100-x}{100}$.
\end{itemize}}

\begin{listedexemples}
    \item \fbox{$\qty{20}{\EURO} + 10~\%$ :} \hspace{2cm} $\qty{20}{\EURO} \times \dfrac{100 + 10}{100} = \qty{20}{\EURO} \times \dfrac{110}{100} = \qty{22}{\EURO}$ 
    \item \fbox{$\qty{150}{\EURO} - 30~\%$ :} \hspace{2cm} $\qty{150}{\EURO} \times \dfrac{100 - 30}{100} = \qty{150}{\EURO} \times \dfrac{70}{100} = \qty{105}{\EURO} $
\end{listedexemples}

\regle{Le pourcentage d'augmentation ou de diminution entre une valeur de départ et une valeur d'arrivée est :}

$$ \mbox{valeur d'arrivée} - \mbox{valeur de départ} $$

\partie{Échelle}

\definition{Dans un plan, l'échelle est le coefficient de proportionnalité tel que : $$\mbox{échelle} = \dfrac{\mbox{distance sur le plan}}{\mbox{distance réelle}}$$}

\remarque{Il faut la même unité pour les deux distances dans le calcul de l'échelle.}

\exemple{Sur un plan, \qty{1}{cm} sur le plan correspond à \qty{50}{km} dans la réalité. $$\mbox{échelle} = \dfrac{\qty{1}{cm}}{\qty{50}{km}}$$}

\propriete{Si les longueurs d'une figure sont multipliées par un facteur $k$, les aires sont multipliées par un facteur $k^2$ et les volumes par un facteur $k^3 = k \times k \times k$.}

\exemple{Si les longueurs sont multipliées par 5, alors les aires sont multipliées par $5^2 = 5 \times 5 = 25$, et les volumes sont multipliés par $5^3 = 5 \times 5 \times 5 = 125$.}

\clearpage
\EXERCICES

\begin{questions}   

    \exercicetitre{Additionner des fractions}
        \newcommand{\exAdditionFraction}[4]{
            $\dfrac{\mathstrut#1}{\mathstrut#2} + \dfrac{\mathstrut#3}{\mathstrut#4} = $
        }
        \begin{multicols}{3}
            \question \exAdditionFraction{20}{4}{2}{4}
            \question \exAdditionFraction{19}{12}{10}{8}
            \question \exAdditionFraction{17}{3}{5}{5}
            \question \exAdditionFraction{20}{5}{6}{14}
            \question \exAdditionFraction{9}{16}{6}{16}
            \question \exAdditionFraction{20}{18}{17}{14}
            \question \exAdditionFraction{15}{7}{7}{8}
            \question \exAdditionFraction{5}{4}{3}{5}
            \question \exAdditionFraction{7}{4}{9}{6}
            \question \exAdditionFraction{16}{19}{14}{12}
            \question \exAdditionFraction{1}{19}{17}{20}
            \question \exAdditionFraction{3}{2}{14}{11}
            \question \exAdditionFraction{18}{19}{7}{18}
            \question \exAdditionFraction{17}{3}{12}{12}
            \question \exAdditionFraction{9}{19}{19}{14}
        \end{multicols}
    
    \exercicetitre{Trouver une fraction égale à une autre}
        \newcommand{\exFractionEgale}[3]{
            $\dfrac{\mathstrut#1}{\mathstrut#2} = \dfrac{\mathstrut#3}{\mathstrut ?}$
        }
        \begin{multicols}{3}
            \question \exFractionEgale{3}{5}{9}
            \question \exFractionEgale{9}{17}{63}
            \question \exFractionEgale{5}{13}{45}
            \question \exFractionEgale{17}{6}{136}
            \question \exFractionEgale{7}{18}{28}
            \question \exFractionEgale{16}{8}{64}
            \question \exFractionEgale{9}{20}{72}
            \question \exFractionEgale{20}{11}{80}
            \question \exFractionEgale{10}{12}{90}
            \question \exFractionEgale{2}{20}{12}
            \question \exFractionEgale{6}{2}{48}
            \question \exFractionEgale{12}{11}{72}
            \question \exFractionEgale{6}{15}{18}
            \question \exFractionEgale{6}{12}{30}
            \question \exFractionEgale{3}{18}{12}
            \question \exFractionEgale{10}{14}{70}
            \question \exFractionEgale{18}{3}{126}
            \question \exFractionEgale{16}{2}{48}
            \question \exFractionEgale{2}{12}{16}
            \question \exFractionEgale{20}{13}{160}
            \question \exFractionEgale{12}{19}{108}
        \end{multicols}

    \exercicetitre{Multiplications de fractions}
        \newcommand{\exMultiplicationFractions}[4]{
            $\dfrac{\mathstrut#1}{\mathstrut#2} \times \dfrac{\mathstrut#3}{\mathstrut#4} = $
        }
        \begin{multicols}{3}
            \question \exMultiplicationFractions{4}{5}{2}{8}
            \question \exMultiplicationFractions{12}{5}{3}{11}
            \question \exMultiplicationFractions{9}{8}{7}{9}
            \question \exMultiplicationFractions{8}{15}{1}{3}
            \question \exMultiplicationFractions{3}{6}{4}{11}
            \question \exMultiplicationFractions{13}{12}{1}{11}
            \question \exMultiplicationFractions{6}{12}{1}{2}
            \question \exMultiplicationFractions{12}{5}{6}{4}
            \question \exMultiplicationFractions{15}{12}{2}{10}
            \question \exMultiplicationFractions{14}{10}{15}{3}
            \question \exMultiplicationFractions{1}{5}{1}{1}
            \question \exMultiplicationFractions{15}{7}{4}{15}
            \question \exMultiplicationFractions{12}{11}{11}{5}
            \question \exMultiplicationFractions{11}{15}{15}{1}
            \question \exMultiplicationFractions{8}{6}{2}{4}
        \end{multicols}

    \clearpage
    \exercicetitre{Divisions de fractions}
        \newcommand{\exDivisionFractions}[4]{
            $\dfrac{\mathstrut#1}{\mathstrut#2} \div \dfrac{\mathstrut#3}{\mathstrut#4} = $
        }
        \begin{multicols}{3}
            \question \exDivisionFractions{3}{9}{2}{9}
            \question \exDivisionFractions{12}{5}{1}{1}
            \question \exDivisionFractions{2}{1}{10}{3}
            \question \exDivisionFractions{15}{10}{12}{15}
            \question \exDivisionFractions{5}{5}{10}{2}
            \question \exDivisionFractions{5}{2}{1}{3}
            \question \exDivisionFractions{1}{2}{1}{14}
            \question \exDivisionFractions{4}{10}{5}{6}
            \question \exDivisionFractions{2}{8}{12}{7}
            \question \exDivisionFractions{8}{6}{12}{10}
            \question \exDivisionFractions{1}{10}{7}{10}
            \question \exDivisionFractions{15}{6}{11}{11}
            \question \exDivisionFractions{6}{12}{9}{1}
            \question \exDivisionFractions{6}{15}{10}{10}
        \end{multicols}
\end{questions}

\clearpage
\CORRECTIONS

\begin{questions}
    \exercicetitre{Additionner des fractions}
        \newcommand{\corAdditionFraction}[4]{%
            % notation : a/b + c/d
            % si les dénominateurs sont égaux
            \ifnumcomp{#2}{=}{#4}{
                $\dfrac{\mathstrut#1}{\mathstrut#2} + \dfrac{\mathstrut#3}{\mathstrut#4} = \dfrac{\pgfmathparse{int(#1+#3)}\pgfmathresult}{#2}$
            }{%
                % si b > d et b % d = 0
                \pgfmathsetmacro{\modulo}{int(mod(#2,#4))}%
                \ifboolexpr{
                    test {\ifnumgreater{#2}{#4}}
                    and
                    test {\ifnumequal{\modulo}{0}}
                }{%
                    $\dfrac{\mathstrut#1}{\mathstrut#2} + \dfrac{\mathstrut#3}{\mathstrut#4} 
                    = \dfrac{\mathstrut#1}{\mathstrut#2} + \dfrac{\mathstrut\pgfmathparse{int(#3*#2/#4)}}{\mathstrut#2} 
                    = \dfrac{\pgfmathparse{int(#1+#3*#2/#4)}\pgfmathresult}{#2}$
                }{}%
                % si d > b et d % b = 0
                \pgfmathsetmacro{\modulo}{int(mod(#4,#2))}%
                \ifboolexpr{
                    test {\ifnumgreater{#4}{#2}}
                    and
                    test {\ifnumequal{\modulo}{0}}
                }{
                    $\dfrac{\mathstrut#1}{\mathstrut#2} + \dfrac{\mathstrut#3}{\mathstrut#4} 
                    = \dfrac{\mathstrut\pgfmathparse{int(#1*#4/#2)}\pgfmathresult}{\mathstrut#4} + \dfrac{\mathstrut#3}{\mathstrut#4} 
                    = \dfrac{\pgfmathparse{int(#1*#4/#2+#3)}\pgfmathresult}{#4}$
                }{}%
                % cas général
                \pgfmathsetmacro{\moduloA}{int(mod(#2,#4))}%
                \pgfmathsetmacro{\moduloB}{int(mod(#4,#2))}%
                \ifboolexpr{
                    not ( test {\ifnumequal{\moduloA}{0}} )
                    and
                    not ( test {\ifnumequal{\moduloB}{0}} )
                }{
                    $\dfrac{\mathstrut#1}{\mathstrut#2} + \dfrac{\mathstrut#3}{\mathstrut#4}
                    = \dfrac{\pgfmathparse{int(#1*#4)}\pgfmathresult}{\pgfmathparse{int(#2*#4)}\pgfmathresult} + \dfrac{\pgfmathparse{int(#2*#3)}\pgfmathresult}{\pgfmathparse{int(#2*#4)}\pgfmathresult}
                    = \dfrac{\pgfmathparse{int(#1*#4+#2*#3)}\pgfmathresult}{\pgfmathparse{int(#2*#4)}\pgfmathresult}$
                }{}%
            }
        }
        \begin{multicols}{2}
            \question \corAdditionFraction{20}{4}{2}{4}
            \question \corAdditionFraction{19}{12}{10}{8}
            \question \corAdditionFraction{17}{3}{5}{5}
            \question \corAdditionFraction{20}{5}{6}{14}
            \question \corAdditionFraction{9}{16}{6}{16}
            \question \corAdditionFraction{20}{18}{17}{14}
            \question \corAdditionFraction{15}{7}{7}{8}
            \question \corAdditionFraction{5}{4}{3}{5}
            \question \corAdditionFraction{7}{4}{9}{6}
            \question \corAdditionFraction{16}{19}{14}{12}
            \question \corAdditionFraction{1}{19}{17}{20}
            \question \corAdditionFraction{3}{2}{14}{11}
            \question \corAdditionFraction{18}{19}{7}{18}
            \question \corAdditionFraction{17}{3}{12}{12}
            \question \corAdditionFraction{9}{19}{19}{14}
        \end{multicols}

    \exercicetitre{Trouver une fraction égale à une autre}
        \newcommand{\corFractionEgale}[3]{
            \begin{tikzpicture}[baseline]
                \node at (0,0) {$\dfrac{\mathstrut#1}{\mathstrut#2} = \dfrac{\mathstrut#3}{\mathstrut \textcolor{rouge}{\pgfmathparse{int(#2*#3/#1)}\pgfmathresult}}$};
                \draw[rouge,-Latex] (-0.45,0.6) arc(180:0:0.5 and 0.15) node[midway,above]{\small $\times\pgfmathparse{int(#3/#1)}\pgfmathresult$};
            \end{tikzpicture}
        }
        \begin{multicols}{3}
            \question \corFractionEgale{3}{5}{9}
            \question \corFractionEgale{9}{17}{63}
            \question \corFractionEgale{5}{13}{45}
            \question \corFractionEgale{17}{6}{136}
            \question \corFractionEgale{7}{18}{28}
            \question \corFractionEgale{16}{8}{64}
            \question \corFractionEgale{9}{20}{72}
            \question \corFractionEgale{20}{11}{80}
            \question \corFractionEgale{10}{12}{90}
            \question \corFractionEgale{2}{20}{12}
            \question \corFractionEgale{6}{2}{48}
            \question \corFractionEgale{12}{11}{72}
            \question \corFractionEgale{6}{15}{18}
            \question \corFractionEgale{6}{12}{30}
            \question \corFractionEgale{3}{18}{12}
            \question \corFractionEgale{10}{14}{70}
            \question \corFractionEgale{18}{3}{126}
            \question \corFractionEgale{16}{2}{48}
            \question \corFractionEgale{2}{12}{16}
            \question \corFractionEgale{20}{13}{160}
            \question \corFractionEgale{12}{19}{108}
        \end{multicols}


    \exercicetitre{Multiplication de fractions}
        \newcommand{\corMultiplicationFractions}[4]{
            $\dfrac{\mathstrut#1}{\mathstrut#2} \times \dfrac{\mathstrut#3}{\mathstrut#4} = \dfrac{#1 \times #3}{#2 \times #4} = \dfrac{\pgfmathparse{int(#1*#3)}\pgfmathresult}{\pgfmathparse{int(#2*#4)}\pgfmathresult}$
        }
        \begin{multicols}{2}
            \question \corMultiplicationFractions{4}{5}{2}{8}
            \question \corMultiplicationFractions{12}{5}{3}{11}
            \question \corMultiplicationFractions{9}{8}{7}{9}
            \question \corMultiplicationFractions{8}{15}{1}{3}
            \question \corMultiplicationFractions{3}{6}{4}{11}
            \question \corMultiplicationFractions{13}{12}{1}{11}
            \question \corMultiplicationFractions{6}{12}{1}{2}
            \question \corMultiplicationFractions{12}{5}{6}{4}
            \question \corMultiplicationFractions{15}{12}{2}{10}
            \question \corMultiplicationFractions{14}{10}{15}{3}
            \question \corMultiplicationFractions{1}{5}{1}{1}
            \question \corMultiplicationFractions{15}{7}{4}{15}
            \question \corMultiplicationFractions{12}{11}{11}{5}
            \question \corMultiplicationFractions{11}{15}{15}{1}
            \question \corMultiplicationFractions{8}{6}{2}{4}
        \end{multicols}

    \exercicetitre{Divisions de fractions}
        \newcommand{\corDivisionFractions}[4]{
            $\dfrac{\mathstrut#1}{\mathstrut#2} \div \dfrac{\mathstrut#3}{\mathstrut#4} = \dfrac{\mathstrut#1}{\mathstrut#2} \times \dfrac{\mathstrut#4}{\mathstrut#3} = \dfrac{#1 \times #4}{#2 \times #3} = \dfrac{\pgfmathparse{int(#1*#4)}\pgfmathresult}{\pgfmathparse{int(#2*#3)}\pgfmathresult}$
        }
        \begin{multicols}{2}
            \questionX \corDivisionFractions{3}{9}{2}{9}
            \questionX \corDivisionFractions{12}{5}{1}{1}
            \questionX \corDivisionFractions{2}{1}{10}{3}
            \questionX \corDivisionFractions{15}{10}{12}{15}
            \questionX \corDivisionFractions{5}{5}{10}{2}
            \questionX \corDivisionFractions{5}{2}{1}{3}
            \questionX \corDivisionFractions{1}{2}{1}{14}
            \questionX \corDivisionFractions{4}{10}{5}{6}
            \questionX \corDivisionFractions{2}{8}{12}{7}
            \questionX \corDivisionFractions{8}{6}{12}{10}
            \questionX \corDivisionFractions{1}{10}{7}{10}
            \questionX \corDivisionFractions{15}{6}{11}{11}
            \questionX \corDivisionFractions{6}{12}{9}{1}
            \questionX \corDivisionFractions{6}{15}{10}{10}
        \end{multicols}

\end{questions}


\end{document}